\section{Introduction}

From 1818 to 1832 Gauss directed the triangulation \index{triangulation} of the Kingdom of Hannover, covering an area of more than 15{,}000 square kilometers. He invented the heliotrope \index{heliotrope} to reflect sunlight so that angles could be measured over distances of 50--100 kilometers, and he developed exactly the mathematics needed to make such a survey reliable: triangle chains built from a single measured baseline, the spherical excess \index{spherical excess} that corrects for the curvature of the Earth, and -- decisively -- the method of least squares, which he invented precisely to reconcile the inconsistencies of redundant angle measurements. Gauss reported the survey in his \textit{Werke}~\cite{gauss-works} and applied the same ideas to astronomy in the \textit{Theoria Motus}~\cite{gauss-theoria}.

In this chapter, we cover:
\begin{itemize}
    \item Great-circle and ellipsoidal distances (haversine and Vincenty \index{Vincenty's formula})
    \item Forward triangulation from a baseline using the sine rule
    \item The three-point (resection) problem \index{three-point problem}
    \item Spherical excess and the area of spherical triangles
    \item Angle adjustment and least-squares adjustment of a triangulation network
\end{itemize}

\section{Measuring Distance on the Earth}

A survey begins with one carefully measured baseline; every other distance is derived by trigonometry. On a sphere of radius $R$, the great-circle distance between two points at (latitude, longitude) coordinates is given by the haversine formula \index{haversine formula}:
\[
a = \sin^2\!\left(\frac{\Delta\varphi}{2}\right) + \cos\varphi_1 \cos\varphi_2 \sin^2\!\left(\frac{\Delta\lambda}{2}\right), \qquad
d = 2R \arctan\!\left(\frac{\sqrt{a}}{\sqrt{1-a}}\right).
\]
Because the Earth is flattened at the poles, Gauss's geodetic work used the ellipsoid; the Vincenty inverse formula computes ellipsoidal distances to sub-millimeter accuracy on the WGS84 ellipsoid.

\section{Triangulation Chains}

The heart of the method is to turn a baseline and two angles into a third point. Suppose the baseline runs from $P_1$ to $P_2$, and at each endpoint the angle between the baseline and the line to the unknown point $P$ is measured. The two rays from $P_1$ and $P_2$ meet at $P$; by the sine rule the sides are
\[
\overline{P_1 P} = c\,\frac{\sin \beta}{\sin(\pi - \alpha - \beta)}, \qquad
\overline{P_2 P} = c\,\frac{\sin \alpha}{\sin(\pi - \alpha - \beta)},
\]
where $c$ is the baseline length. Repeating this triangle-by-triangle, the coordinates of an entire chain of points follow from one baseline alone.

\begin{example}[The First Triangle]
Given the baseline $P_1 = (0,0)$, $P_2 = (1,0)$ and angles $\alpha = \beta = 60^\circ$, the module computes the apex $(0.5, \sqrt{3}/2)$ and both derived sides exactly 1:
\begin{lstlisting}[style=python]
>>> from geodesy import forward_triangulation
>>> from math import radians
>>> pt, d1, d2 = forward_triangulation((0, 0), (1, 0), radians(60), radians(60))
>>> print(f"apex = ({pt[0]:.4f}, {pt[1]:.4f}), sides = {d1:.4f}, {d2:.4f}")
apex = (0.5000, 0.8660), sides = 1.0000, 1.0000
\end{lstlisting}
\end{example}

\section{The Three-Point Problem}

When the observer stands at the unknown point $P$ and measures the angles $\alpha = \angle APB$ and $\beta = \angle BPC$ between three known stations $A, B, C$, the position of $P$ is determined by intersecting the two circles on which $AB$ subtends angle $\alpha$ and $BC$ subtends angle $\beta$. This is the classic \index{three-point problem} Snellius--Pothenot problem, familiar to every surveyor (Gauss used it constantly).

\begin{proposition}
The circle through $A$ and $B$ from which the chord $AB$ subtends angle $\alpha$ has radius $c/(2\sin\alpha)$ and center at the midpoint of $AB$ shifted a distance $(c/2)\cot\alpha$ along the perpendicular, where $c = \lvert AB \rvert$.
\end{proposition}

The point $P$ is the second intersection of the two such circles (one for $AB$, one for $BC$), and it is located by verifying that the computed angles match the observed ones.

\begin{example}[Resection]
\begin{lstlisting}[style=python]
>>> from geodesy import resection
>>> A, B, C = (0.0, 0.0), (1.0, 0.0), (2.0, 1.0)
>>> # angles observed from the unknown station P:
>>> alpha, beta = radians(56.3099), radians(75.9638)
>>> resection(A, B, C, alpha, beta)
(0.7, 0.9)
\end{lstlisting}
The measured angles pinpoint the observer to within the precision of the circle intersection.
\end{example}

\section{Spherical Excess}

A triangle on the Earth's surface is not Euclidean: its angles sum to more than $\pi$, and the excess measures its area. Euler and L'Huilier gave the formula used in surveying:
\[
\tan\frac{E}{4} = \sqrt{\tan\frac{s}{2}\tan\frac{s-a}{2}\tan\frac{s-b}{2}\tan\frac{s-c}{2}},
\]
where $a, b, c$ are the spherical side lengths (in radians) and $s = (a+b+c)/2$. The excess $E$ is the area on the unit sphere, so the true area is $E R^2$.

\begin{example}[The Octant]
\begin{lstlisting}[style=python]
>>> from geodesy import spherical_excess
>>> e, area = spherical_excess(0, 0, 0, 90, 90, 0)
>>> print(f"excess = {e:.4f} rad = {e * 180 / 3.14159:.1f} deg")
>>> print(f"area   = {area:.2e} m^2")
excess = 1.5708 rad = 90.0 deg
area   = 6.38e+13 m^2
\end{lstlisting}
The octant triangle $(0,0)$, $(0,90^\circ)$, $(90^\circ,0)$ covers one eighth of the sphere, so its area is $4\pi R^2/8 = \pi R^2/2$, and the excess is the full $\pi/2$.
\end{example}

For the small triangles of a national survey the excess is seconds of arc; Gauss's survey of Hannover nonetheless computed it at every triangle to close the angle sums on the ellipsoid.

\section{Adjustment of a Network}

Every measurement carries error, so a network of triangles never closes exactly. Gauss's solution is the method of least squares \index{method of least squares}: distribute the smallest possible corrections (weighted by their variances) subject to the geometric constraints. For angle observations this is a \emph{conditional} adjustment: minimize $c^T P c$ subject to $A c = -w$, where $c$ is the vector of corrections, $w$ the misclosures (the amount by which angle sums miss $180^\circ$), and $A$ encodes which angles belong to which triangle. The solution is
\[
c = -P^{-1} A^T (A P^{-1} A^T)^{-1} w,
\]
the same normal-equation structure Gauss used everywhere from astronomy~\cite{gauss-theoria} to surveying; see \cite{bjorck} for a modern treatment.

\begin{example}[Closing Two Triangles]
Suppose six observed angles form two triangles with sums $180.002^\circ$ and $180.001^\circ$. Equal-weight adjustment distributes each misclosure equally:
\begin{lstlisting}[style=python]
>>> import numpy as np
>>> from geodesy import conditional_adjustment
>>> obs = [60.001, 60.002, 59.999, 60.003, 59.998, 60.000]
>>> w = np.array([sum(obs[:3]) - 180.0, sum(obs[3:]) - 180.0])
>>> A = np.zeros((2, 6)); A[0, :3] = 1; A[1, 3:] = 1
>>> c = conditional_adjustment(A, w)
>>> print(f"misclosures = {w}, corrections = {[round(x, 5) for x in c]}")
misclosures = [0.002 0.001], corrections = [-0.00067, -0.00067, -0.00067, -0.00033, -0.00033, -0.00033]
\end{lstlisting}
The corrected sums both close at $180.000^\circ$, and the three angles of each triangle share the correction equally, as Gauss's corner adjustment prescribes.
\end{example}

\section{Python Implementation}

In \texttt{python/geodesy.py}:

\begin{lstlisting}[style=python, caption={Key functions in python/geodesy.py}]
haversine_distance(lat1, lon1, lat2, lon2)   # Spherical distance (m)
vincenty_distance(lat1, lon1, lat2, lon2)    # Ellipsoidal WGS84 distance (m)
spherical_excess(lat1, lon1, lat2, lon2, lat3, lon3)
forward_triangulation(p1, p2, angle1, angle2)  # Baseline -> third vertex
intersection(p1, p2, alpha1, alpha2)           # Two-ray intersection
resection(p1, p2, p3, alpha, beta)             # Three-point problem
adjust_triangle_angles(angles)                 # Equal-correction closing
conditional_adjustment(A, w, weights)          # Least-squares adjustment
meridian_arc_degree(lat)                       # 1 degree of latitude (m)
\end{lstlisting}

\section{Numerical Examples}

\begin{example}[Meridian Arc]
The length of one degree of latitude increases from 110.6 km at the equator to 111.7 km near the pole -- the geometric signature of the ellipsoid that Gauss's survey had to account for:
\begin{lstlisting}[style=python]
>>> from geodesy import meridian_arc_degree
>>> meridian_arc_degree(0)
110574.4
>>> meridian_arc_degree(45)
111141.5
\end{lstlisting}
\end{example}

\begin{example}[Ellipsoidal Distance]
\begin{lstlisting}[style=python]
>>> from geodesy import vincenty_distance
>>> vincenty_distance(48.8566, 2.3522, 51.5074, -0.1278)
343923.1
\end{lstlisting}
Paris to London, measured on the WGS84 ellipsoid: 343.9 km. The same formula, applied to the chain of triangles, is how the Hannover survey derived every distance in the kingdom from a few measured baselines.
\end{example}

\begin{example}[Triangle Adjustment]
\begin{lstlisting}[style=python]
>>> from geodesy import adjust_triangle_angles
>>> adjust_triangle_angles([59.95, 60.05, 60.03])
[59.94, 60.04, 60.02]
\end{lstlisting}
The misclosure of $0.03^\circ$ is divided equally among the three corners.
\end{example}

\section{Exercises}

\begin{exercise}
Verify that \texttt{forward\_triangulation} on the baseline $(0,0)$--$(1,0)$ with angles $60^\circ$, $60^\circ$ returns the apex of an equilateral triangle, and compute the third angle of the triangle.
\end{exercise}

\begin{exercise}
A baseline of length 2 km runs east--west. Angles measured from its ends to a hilltop are $45^\circ$ and $75^\circ$. Find the distances to the hilltop by the sine rule and by \texttt{forward\_triangulation}.
\end{exercise}

\begin{exercise}
Show that the sphere octant $(0,0)$, $(0,90^\circ)$, $(90^\circ,0)$ has spherical excess $90^\circ$, and verify with \texttt{spherical\_excess}. Why can a terrestrial triangle never have excess exceeding $90^\circ$?
\end{exercise}

\begin{exercise}
Use \texttt{resection} to locate a point $P$ given stations $A = (0,0)$, $B = (1,0)$, $C = (2,1)$ and measured angles $\alpha = 56.31^\circ$, $\beta = 75.96^\circ$. Check the result by recomputing the angles.
\end{exercise}

\begin{exercise}
Two triangles in a network have misclosures of $0.005^\circ$ and $0.010^\circ$. Set up the conditional adjustment and show the corrections with equal weights, then with the second triangle's angles weighted half as much.
\end{exercise}

\begin{exercise}
Explain why the length of a degree of latitude grows toward the pole, and compute \texttt{meridian\_arc\_degree} at the equator and at $80^\circ$ to quantify the difference.
\end{exercise}

\section{Exercise Solutions}

\begin{solution}
The third vertex is $(0.5, \sqrt{3}/2)$ and the derived sides are both exactly 1, so the triangle is equilateral. The third angle is $180^\circ - 60^\circ - 60^\circ = 60^\circ$.
\end{solution}

\begin{solution}
With baseline $c = 2$ km and angles $\alpha = 45^\circ$, $\beta = 75^\circ$, the third angle is $60^\circ$. By the sine rule, $\overline{P_1P} = c \sin 75^\circ/\sin 60^\circ \approx 2.232$ km and $\overline{P_2P} = c \sin 45^\circ/\sin 60^\circ \approx 1.633$ km; \texttt{forward\_triangulation} returns the same values.
\end{solution}

\begin{solution}
All three sides of the octant are quarter great circles, so $a = b = c = \pi/2$ and $s = 3\pi/4$. L'Huilier's formula gives $\tan(E/4) = \tan(\pi/8)$, hence $E = \pi/2 = 90^\circ$. The area is $4\pi R^2/8$, so $E = \text{area}/R^2 = \pi/2$. A spherical triangle cannot exceed a hemisphere, whose excess is $2\pi$ steradians, i.e.\ $90^\circ$ per octant bound; in any event the excess of a survey triangle is seconds of arc.
\end{solution}

\begin{solution}
The module returns $(0.7, 0.9)$; recomputing the angles with the law of cosines reproduces $56.31^\circ$ and $75.96^\circ$ to the precision of the printed values, confirming the point.
\end{solution}

\begin{solution}
For two triangles with six angles, $A$ is the $2 \times 6$ matrix with ones in the first three and last three columns, and $w = (0.005, 0.010)$. Equal weights give corrections $-0.00167$ for each angle of the first triangle and $-0.00333$ for each of the second. Halving the weights of the second triangle's angles doubles their corrections to $-0.00667$, concentrating the adjustment where the measurements are less reliable.
\end{solution}

\begin{solution}
The meridional radius of curvature of the ellipsoid increases with latitude (the flattening shortens the polar radius), so one degree of latitude is longer near the pole. The module gives $110574.4$ m at the equator and $111663.2$ m at $80^\circ$ -- a difference of about 1.1 km, or roughly 1\%.
\end{solution}

\section{References}

\begin{itemize}
    \item \cite{gauss-works} -- Gauss, \textit{Werke}, Vols.~IV, VIII--IX (Hannover survey, geodetic papers)
    \item \cite{gauss-theoria} -- Gauss, \textit{Theoria Motus} (1809), the least-squares method
    \item \cite{bjorck} -- Bj\"orck, \textit{Numerical Methods for Least Squares Problems}
\end{itemize}
